% Requires in preamble: \usepackage{booktabs, amsmath, cleveref}
\subsection{Renovation Strategies and 2050 Targets}
Since energy performance is a proxy for a building's energy efficiency, improving
it is a central policy objective pursued through several frameworks and
strategies. At the European level, the Energy Performance of Buildings Directive
(EPBD) is the primary legal instrument for decarbonising the building stock by
2050: it mandates progressive renovation of existing buildings, zero-emission
standards for new construction, and a phase-out of fossil-fuel heating, while
promoting on-site renewable generation \cite{epbd_2024}. In Belgium, however, the
three regions differ in both the strategy and the methodology used to translate
these objectives into national targets, as summarised in
\Cref{tab:regional_renovation_targets}.

\begin{table}[ht]
  \centering
  \caption{Regional 2050 residential renovation targets and policy-compatible
           renovation rates.}
  \label{tab:regional_renovation_targets}
  \begin{tabular}{lllcl}
    \toprule
    Region & Label system & 2050 average target & Energy score & Renovation rate \\
           &              &                     & (kWh/m$^2\cdot$yr) & (policy-compatible) \\
    \midrule
    Flanders & EPC & label A                    & 100 & $\sim$3\%/yr \\
    Wallonia & PEB & decarbonised label A        & 85  & $>$3\%/yr \\
    Brussels & PEB & C+                          & 100 & $\sim$2.5\%/yr\textsuperscript{a} \\
    \bottomrule
  \end{tabular}

  \vspace{2pt}
  {\footnotesize \textsuperscript{a} RENOLUTION states that the current
   renovation pace is insufficient but does not fix an official annual rate;
   $\sim$2.5\%/yr is the policy-compatible assumption adopted here, consistent
   with the Belgium-wide 2.5--3\%/yr requirement.}
\end{table}

\paragraph{Flemish Region.} Flanders defines its 2050 target sharply: existing
residential buildings should reach a performance level comparable to a new
dwelling built under a 2015 permit, i.e.\ EPC label~A with an energy score of
100~kWh/m$^2$. Reaching this level requires renovating, on average, about 3\% of
existing dwellings per year \cite{veka_ltrs_2050}.

\paragraph{Walloon Region.} Wallonia expresses its target less as a single
kWh/m$^2$ threshold and more as a stock-level objective: by 2050 the residential
stock should tend towards an average \emph{decarbonised} PEB label~A (the upper
threshold of which is 85~kWh/m$^2$), with priority given to the deep renovation of
the worst-performing dwellings. The strategy likewise notes that meeting this
objective requires raising the energy-renovation rate above 3\%/year
\cite{wallonia_ltrs_2020}.

\paragraph{Brussels-Capital Region.} The RENOLUTION strategy sets a clear
numerical 2050 target: the average residential PEB level should reach C+,
corresponding to 100~kWh/m$^2\cdot$year, roughly a threefold reduction relative
to the current average. It also defines dwelling-level milestones, requiring every
dwelling to reach at least PEB~275 by 2033 and, later, PEB~150
\cite{renolution_brussels}.

For the renovation-rate assumptions, the Belgium-wide climate-neutrality
scenarios indicate that the building-sector renovation rate must rise from about
1\%/year today to 2.5--3\%/year, with the mix shifting from shallow to
predominantly deep renovations \cite{climat_neutral_be_2050}. These regional 2050
targets and renovation-rate ranges provide the policy basis for the renovation
scenarios applied to the archetypes of \Cref{sec:size_compo}.
